% Generated by maint/texbuild/theory_tex.py from theory/workbooks/cell_tracking/build_plan.md.
% Edits here are overwritten. Edit the markdown, then run the script again.
\chapter{\texorpdfstring{The build plan: the cell program}{The build plan: the cell program}}
\label{chap:build-plan}

\textbf{Purpose:} The order the cell program's wants in \hyperref[chap:cell-tracking-table]{cell\_\allowbreak{}tracking\_\allowbreak{}table.md} are built in, and every ruling on them, dated, so every session builds the next item and not whatever looks nearest. Built in order, one item at a time; an item is done when its program runs on the device and its row below says so.

\textbf{Two tables, two concerns (Doug, 23 September):} keep the cell table and the engine table separate; each is optimized against the other only through its interface section, never merged.

\textbf{Scope:} every want in the cell table's “wants” column. Each one is a UTM program: a record program swept over the body magnitudes, or a linear key swept over the voxels. Each has its own program module (like \texttt{residual\_\allowbreak{}program} and \texttt{fingerprint\_\allowbreak{}program}), its own \texttt{-\allowbreak{}-\allowbreak{}run} part, and its own measurement. Ingest stays its own step and is never part of a run. The item numbers are the plan's own, kept from 22 to 24 September, so a number missing here is the engine book's (build\_\allowbreak{}plan.md).

The bodies are one array: the tracker's leaves are the \texttt{max\_\allowbreak{}tree} bodies in the order \texttt{flatten} packs them, so a body's index is its frame's start plus its leaf.

Statuses: proved, measured, built, theory, refuted, not so. “Not built” means nothing of it exists yet.

\section{\texorpdfstring{Done before this plan}{Done before this plan}}

{\small
\setlength{\tabcolsep}{6pt}
\begin{longtable}{>{\RaggedRight\arraybackslash}p{\dimexpr0.1639\linewidth-2\tabcolsep\relax}>{\RaggedRight\arraybackslash}p{\dimexpr0.4573\linewidth-2\tabcolsep\relax}>{\RaggedRight\arraybackslash}p{\dimexpr0.3788\linewidth-2\tabcolsep\relax}}
\toprule
item & what it is & status \\
\midrule
\endhead
two steps & \texttt{-\allowbreak{}-\allowbreak{}ingest} alone; \texttt{-\allowbreak{}-\allowbreak{}run <part>} repeated, in order & built \\
rows 2, 15: the print & \texttt{fingerprint\_\allowbreak{}program}, 52 steps: the mass band and six signed golden bands of K′/m², K′ = sₐs\_\allowbreak{}b(m·M − SₐS\_\allowbreak{}b), scales \texttt{voxel\_\allowbreak{}pm}/gcd & built: 81,807 bodies of 44b6\_\allowbreak{}0113de3b in 11.7 ms, CRC-64 \texttt{590f3a2c386647e3} twice \\
\bottomrule
\end{longtable}
}

The whole tracker is written out as one equation, stage by stage, in the engine table's “The whole equation” (Doug, 22 September: the aim is the entire equation, so algebra can improve the whole). Every item below that changes a stage updates that section too.

\section{\texorpdfstring{Phase B: wants over the body magnitudes (record programs)}{Phase B: wants over the body magnitudes (record programs)}}

\tablerecord{B1}
\begin{description}
\item[{row}] 7
\item[{what it is}] a pair program gives each candidate's print cost (overlap triples and the climbed forward); \texttt{heaviest\_\allowbreak{}matching} picks one to one and rewrites the forwards
\item[{status}] built (\texttt{-\allowbreak{}-\allowbreak{}print-\allowbreak{}match}, \texttt{print\_\allowbreak{}pair}, FIELD\_\allowbreak{}SIGNED added to read records back): a body against itself weighs the ceiling 1,275 on all 81,807, device equals host; in the tracker on 44b6\_\allowbreak{}0113de3b 56,927 pairs swept at once in 217 us, 31,130 of 80,697 leaves matched, 19,113 forwards changed; edges unchanged at 38/12 of 50, because base.cfg's \texttt{unbound} linker picks by overlap weight and centre cost and the forward is only one more candidate in its pool
\end{description}

\tablerecord{B2}
\begin{description}
\item[{row}] 5
\item[{what it is}] a pair program gives the centroid change as an exact rational; a three-body program predicts the landing by cross-multiplying; a tally counts agreement with the climb
\item[{status}] built (\texttt{-\allowbreak{}-\allowbreak{}velocity}, \texttt{velocity}, 42 steps, three members: two bodies and the climbed lag as its own array). The record is the velocity itself, the change Sₐ′m − Sₐm′ over m·m′ per axis, and a truthy verdict per axis: the climbed lag stands on a lattice point beside the dipole's change, |Δ − L| < 1, cross-multiplied as |R| < m·m′, written COMPARE·(COMPARE + 1). On 44b6\_\allowbreak{}0113de3b: the check (each body and its climbed forward) 35,755 lanes, all three axes agree on 8,720 (z 22,660, y 13,227, x 13,306); the prediction (the last change against the next climb) 35,966 lanes, 752 agree (z 13,524, y 5,520, x 5,163); 0 differ from the host in both; edges unchanged at 38/12 of 50
\end{description}

\tablerecord{B3}
\begin{description}
\item[{row}] 6
\item[{what it is}] Doug, 22 September: the membrane (4 nm) is a boundary that deforms, so a body and its match overlap only in part, and by how much is what to calculate. What stays inside the match under every jitter between the frames is the body's internal constituents; their shift is the body's velocity and vector; what falls out is the membrane deforming. Built with C1, which is the same measurement. And (Doug, 22 September) two membranes cannot pass through one another: for every touching pair A, B at t and their matches A′, B′ at t+1, the centroid difference keeps its side, (c\_\allowbreak{}A − c\_\allowbreak{}B)·(c\_\allowbreak{}A′ − c\_\allowbreak{}B′) > 0. Every mass denominator is positive, so the sign is exact from the cross-multiplied numerators: one pass writes each pair's difference as a record, and a second pass reads the two records and writes the truthy verdict. It is cheap here and needs no model
\item[{status}] the deformation folded into C1. The no-crossing pair is built (\texttt{-\allowbreak{}-\allowbreak{}contact-\allowbreak{}side}, \texttt{contact\_\allowbreak{}side}). Pass 1, 23 steps, two bodies: N = σ∘(S\_\allowbreak{}A m\_\allowbreak{}B − S\_\allowbreak{}B m\_\allowbreak{}A), σ = voxel\_\allowbreak{}pm over its common unit. Pass 2, 18 steps, reads two pass-1 records: dot = N·N′, c = COMPARE(dot + 1, 1), kept = c(c + 1), crossed = c(c − 1). The lanes are every touching leaf pair (\texttt{joined}) whose forwards both exist. Each pass is proved against the host. Ran on 44b6\_\allowbreak{}0113de3b and three 6bba: 0 touching pairs on every sample. Today's bodies are separate components of the cut, so none share a face and \texttt{grow\_\allowbreak{}leaves} writes \texttt{joined\_\allowbreak{}count} = 0; E7 (\texttt{group\_\allowbreak{}objects}) is inert too. What “touching” means for components of a cut is waiting on Doug (item 6)
\end{description}

\tablerecord{B4}
\begin{description}
\item[{row}] 9
\item[{what it is}] a three-body program (parent, child, sibling): band(m\_\allowbreak{}c + m\_\allowbreak{}s) = band(m\_\allowbreak{}p), then 3ΔᵀKΔ ≥ tr(K)·|Δ|²
\item[{status}] built (\texttt{-\allowbreak{}-\allowbreak{}division}, \texttt{division}, 100 steps, three members read through an index): the triples are every parent with two pieces whose backward climbs land on it. K is the parent's m·M − S·Sᵀ scaled by voxel\_\allowbreak{}pm over its common unit, and Δ is the children's cross-multiplied centroid difference; both sides are homogeneous, so no mass is divided. On 44b6\_\allowbreak{}0113de3b: 142,649 triples in 13.8 ms; mass conserved on 3,356, split along the parent's long axis on 99,155, both on 2,452; 0 differ from the host. That sample's key holds no division. On 6bba\_\allowbreak{}48816121, 6bba\_\allowbreak{}09961292 and 6bba\_\allowbreak{}cdcfe533: 326,836, 85,146 and 99,724 triples, 0 differ from the host. The keys hold 13 divisions; 3 put the parent and both children on three separate leaves, and 2 of those are among the triples. Mass holds on 1, the axis on 2, both on 1. Why the other 10 are lost (measured, 23 September): both children are on one leaf. The cut at ℓ* has not separated just-divided daughters; no node is off a leaf, and none falls outside consecutive tracked frames. So the division term's loss is row 1's cut, not the climb
\end{description}

\tablerecord{B5}
\begin{description}
\item[{row}] 14
\item[{what it is}] wherever two sizes are compared, the mass band is compared instead (Doug, 22 September: size comparisons only; exact sums stay counts, since a band there would round)
\item[{status}] built (\texttt{-\allowbreak{}-\allowbreak{}mass-\allowbreak{}band}, \texttt{band\_\allowbreak{}or\_\allowbreak{}count} in golden\_\allowbreak{}bands): \texttt{group\_\allowbreak{}objects} (left/right bigger, leans), \texttt{arm\_\allowbreak{}landing}, \texttt{focus\_\allowbreak{}links}, \texttt{assign\_\allowbreak{}bodies} (heaviest) and \texttt{track\_\allowbreak{}bend} (both widest loops) compare bands when the flag is on and counts when it is off. Built clean; A/B not run yet
\end{description}

\tablerecord{B6}
\begin{description}
\item[{row}] 7, 8
\item[{what it is}] OrganoidTracker 2.0's estimation made exact: each link's context-aware probability P(A|G) = Z\_\allowbreak{}A / Z, the sum over every flow-conserving arrangement of the local subgraph (three steps out, divisions allowed) of the product of its weights, with and without link A. Theirs is exp(−E/T) of calibrated network energies. Ours (Doug, 22 September: yes): a link weighs the fraction of the body that lands on the candidate, and “no link” the null draw's fraction; each body takes one outcome in every arrangement, so the mass denominators cancel, and Z\_\allowbreak{}A and Z are exact integers, compared by cross-multiplying, with no temperature and no calibration
\item[{status}] built (\texttt{-\allowbreak{}-\allowbreak{}marginal}, needs the null: \texttt{-\allowbreak{}-\allowbreak{}null} with a floor; \texttt{marginal} module, device kernel one thread a body, host proof on anchor\_\allowbreak{}sift's exact integer, 16 limbs). The local set G is the body, every source that shares one of its candidates, and their candidates, at one frame step. Each source takes one outcome and no two sources take the same target. A body whose set is too wide (more than 15 sources, 64 targets or 2¹⁶ arrangements) is asked alone. \textbf{Differs from the text above:} one step, not three, and no division arrangements (a source to two targets). Those are Doug's to decide. The tally compares the most probable in context, the heaviest candidate alone, and the climb's forward against the key edges. Measured on 44b6\_\allowbreak{}0113de3b: 34,297 bodies asked in 50 ms (30,444 in context, 3,853 alone, 8 too wide), 123,538 options, device equal to the host. Of 49 key edges, the most probable in context is the truth for 16 and “no link” for 27. The heaviest candidate alone gets 38 and the climb's forward 35. The null's weight (the climbed null count) dominates T (counted at the drift); putting them on one footing is waiting on Doug (item 7)
\end{description}

\section{\texorpdfstring{Phase C: wants over the voxels and the history (linear keys, then record gathers)}{Phase C: wants over the voxels and the history (linear keys, then record gathers)}}

\tablerecord{C1}
\begin{description}
\item[{row}] 13, 6
\item[{what it is}] the jitter scrub, with B3 (Doug, 22 September: every width, kept per width, and the jitter between frames). For each body and its match at the climbed lag, a voxel stays at width w when every offset in the box of half-width w around the lag lands it inside the match. That is the match's Chebyshev distance from its own boundary, one exact integer per voxel, so one pass buckets every voxel by the widest box it survives. Mass, sums and moments are re-summed per width: the core that stays is the internal constituents, its shift is the velocity, and what falls out at each width is the membrane deforming
\item[{status}] built 23 September (\texttt{climb\_\allowbreak{}machine\_\allowbreak{}core}, \texttt{-\allowbreak{}-\allowbreak{}core}), proved on 44b6\_\allowbreak{}0113de3b: the width-0 core mass is C at the climbed lag on 35,755 of 35,755 forward climbers. Widths reach 5, and the per-width counts, masses and displacements are in row 13. The deeper the core, the less often its displacement lies within half a voxel of the climbed lag (62\% at w=0, 42\%, 30\%, 13\%, 13\%). Still open: the core velocity feeding a link or the print, and contact κ\_\allowbreak{}w (item 6) built on the same D
\end{description}

\tablerecord{C2}
\begin{description}
\item[{row}] 9, 10
\item[{what it is}] a running-sum key over the \texttt{.\allowbreak{}oapx} flips, then 8 corners gathered per body box: a bump that rises and settles is a division, one that rises and stays is a lysis
\item[{status}] the measurement is built (\texttt{-\allowbreak{}-\allowbreak{}box-\allowbreak{}history}, \texttt{box\_\allowbreak{}history} module, \texttt{climb\_\allowbreak{}machine\_\allowbreak{}extents}). The box extents are held per body from the stored runs; per window and bit, one running-sum key is built on the device, and each body's box is read by its 8 corners. It is proved against walking every box: 0 of 11,780,208 differ on 44b6\_\allowbreak{}0113de3b, 0 of 9,021,600 on 6bba\_\allowbreak{}48816121 (its \texttt{.\allowbreak{}oapx} is written into the scratch set, not the data). Found on 6bba\_\allowbreak{}48816121: every key division's parent is on a leaf of 0.65 to 1.0 M voxels whose box is the whole view, and 881 of 931 key nodes share a leaf with another key node (up to 13 on one leaf). So the bump cannot be read until row 1's cut holds one cell a leaf. The verdict is not built: it waits on the cut and on Doug (Waiting 14)
\end{description}

\tablerecord{C3}
\begin{description}
\item[{row}] 4
\item[{what it is}] a tuple sum over contacts gives each body's departure from the drift, weighted by mass; an origin program traces the flow back toward division and lysis events
\item[{status}] not built
\end{description}

\tablerecord{C4}
\begin{description}
\item[{row}] 12
\item[{what it is}] every departure from the floor placed; the noise keys stamped in one cycle
\item[{status}] not built; waits on Doug's design
\end{description}

C5, the residual as unit sweeps, is the engine's (the engine book's build plan).

\section{\texorpdfstring{Phase D: prove and measure}{Phase D: prove and measure}}

\tablerecord{D3}
\begin{description}
\item[{what it is}] on the 25: time against the 15,362 ms \texttt{-\allowbreak{}-\allowbreak{}run track} baseline, score each want A/B, measure \texttt{above\_\allowbreak{}null} (row 8), and write it all into the ledger and the engine table
\item[{status}] \texttt{-\allowbreak{}-\allowbreak{}mass-\allowbreak{}band} run 23 September: identical to counts on all 24 that ran both sides (6,134 edges, 5,176/227/2/729). It is inert by construction: every reader is gated off on the scored path (see row 14). Timing is not comparable, since the bands side shared the GPU with the core run. 44b6\_\allowbreak{}551a5dba failed its .iapx CRC in the bands run after loading in the counts run, and fails on every read since, with mtime unchanged (Waiting on Doug 13). \texttt{above\_\allowbreak{}null}, measured 23 September on the two samples that have a floor (44b6\_\allowbreak{}0113de3b in the set, 6bba\_\allowbreak{}48816121 in the scratch set): 0.039606 against 0.027942 for every node on 44b6\_\allowbreak{}0113de3b, and 0.001149 against 0.002323 on 6bba\_\allowbreak{}48816121. \texttt{stands} scores 0 on both. The scorer (\texttt{maint/\allowbreak{}score\_\allowbreak{}submission.\allowbreak{}py}, position within 7 µm) finds 2 of 44b6\_\allowbreak{}0113de3b's 50 key edges, where the tracker's leaf-identity tally says 38. The gap, measured edge by edge (scratchpad \texttt{edge\_\allowbreak{}gap.\allowbreak{}py}): of the 100 key-node ends, 11 have a predicted centroid within 7 µm. The rest sit 7 to 15 µm from the nearest centroid, because the leaf holding the key cell is a multi-cell blob (6,516 to 24,481 voxels against about 3,000 for one cell) or a small piece (3 to 357 voxels). On 14 of 50 edges the nearest node's forward is the successor's nearest node. So the leaf-identity tally counts leaf-to-leaf links, while the position scorer cannot find the cell: the loss is row 1's cut again, not the linker. On the 25, \texttt{above\_\allowbreak{}null} needs a floor for every sample, and only one \texttt{.\allowbreak{}oapx} exists in the set. \texttt{-\allowbreak{}-\allowbreak{}print-\allowbreak{}match} (B1) changes forwards, but the linker does not read them (38/12 unchanged). \texttt{-\allowbreak{}-\allowbreak{}unit-\allowbreak{}sweep} is proved identical. So no other built want changes a link on its own
\end{description}

D1 and D2 are the engine's (the engine book's build plan).

\section{\texorpdfstring{Waiting on Doug}{Waiting on Doug}}

\begin{enumerate}
\item Answered 22 September: \texttt{membrane\_\allowbreak{}pm} is 4 nm; B3 is the membrane's deformation, built with C1.
\item The design for placing the floor departures and stamping the noise keys. (C4)
\item B1 rewrites the forwards as the plan says, but the linker base.cfg runs (\texttt{unbound}) does not choose by the forward, so the print cannot move a link there. Should the print weigh in the linker's pool itself (row 7's “link each body to the one whose fingerprint agrees”)?
\item Answered 22 September: B6's weights are the landing fraction and the null draw's fraction.
\item Ultrack (Nature Methods 2025, the Royer lab's own tracker, the one this competition builds on) selects cells from a hierarchy of candidate segments jointly with the links: an ILP keeps one node of every nested chain (y\_\allowbreak{}i + y\_\allowbreak{}j ≤ 1 for i ⊂ j) and maximizes Σ IoU\textasciicircum{}4 over the links. Our component tree is that hierarchy already, and row 1's “one node per cell” is that selection. Their γ = 4 and event penalties are chosen numbers; an exact form would have to replace them.
\item Answered 23 September, touching: “we have assigned a boundary to these things and can use the overlap to estimate contact probability truthy falsy”. Each body has an assigned boundary; two bodies are in contact where their boundaries overlap, and the overlap's size is the truthy magnitude. Built on C1's per-voxel distance: at width w, contact(A, B) = |\{x : d\_\allowbreak{}A(x) ≤ w ∧ d\_\allowbreak{}B(x) ≤ w\}|, where d is the Chebyshev distance to the body. Every width accumulates (item 8).
\item Answered 23 September, the null's weight: “we need the whole equation written out to figure out the exact transform of this term”. From the equation: the null is the climb's own functional A\_\allowbreak{}b(v*) taken in a null frame, and T is a different functional (one target, at the drift). The exact transform splits the same functional by target: T*(b, b′) = |\{x ∈ b : L\_\allowbreak{}\{t+1\}(x + v*\_\allowbreak{}b) = b′\}|, so Σ\_\allowbreak{}\{b′\} T*(b, b′) = h(b). The links and “no link” then come from one operator: the real frame split by target, and the null frame whole.
\item Answered 23 September, C1's units: “everything should accumulate, everything should use the same atom class with their int cast into it”. Every width's sums include every wider width, a running sum over the widths. Every quantity is carried as integers cast into the same atom lanes the engine sweeps, so the box is on the atom's integer lattice, in voxels.
\item Answered 23 September, the cut: “It should be able to slide all the way to the end looking for coherence”. Row 1's selection slides through every level of the component tree, to the top, and keeps what stays coherent between frames, rather than cutting the frame at one level. That is how a just-divided pair and everything else can each be held at their own level.
\item B6's one frame step and the missing division arrangements (see B6).
\item The window of the one operator C: the climb's reach around the drift, or the whole view. Built, 23 September: C at the 27 shifts around each climbed lag plus the drift (every read and proof needs only these). All three proofs hold on 44b6\_\allowbreak{}0113de3b (held count on 726,273 climbers, the local maximum on all of them, the overlap triples on 80,697 bodies). A wider window is only more cells.
\item “No link” in the marginal weighs the best null draw, max\_\allowbreak{}k h⁰\_\allowbreak{}k. With T* (C at v*) the truth is picked for 19 of 49 key edges and “no link” wins 24. Should the draws enter as the best draw, as each draw its own “no link” outcome, or otherwise?
\item \texttt{train/\allowbreak{}44b6\_\allowbreak{}551a5dba/\allowbreak{}44b6\_\allowbreak{}551a5dba.\allowbreak{}iapx} loaded and proved in D3's counts run (about 00:52, 23 September), then failed its CRC in the bands run (about 01:05) and on every read since: rebuilt 6e6947161fd5c5cd against the file's 3245a85b0a90111e. The mtime is still 22 September 19:29. Three MD5 reads since agree (e0c60eb2e06aa3eb618aec5b59599163). The data was not touched from here. Re-ingest it, or check the drive?
\item C2's verdict: which of the 16 bits enter the bump (each separately, or the density Σ count·2\textasciicircum{}bit that the entropy cloud already uses), and how the 11-transition windows meet a one-frame division (the window holding the transition, against its neighbours?). Whatever the answer, it is only readable once the cut (item 9, the sliding cut) puts one cell on a leaf: on 6bba the dividing parents sit on leaves the size of the view. 9b. Answered 23 September, the sliding cut: measure first. Slide ℓ through every level and report, per level, the component count and how many key cells share a leaf; no selection is built until the measurement is read. \textbf{Measured 23 September} (\texttt{-\allowbreak{}-\allowbreak{}slide}, \texttt{max\_\allowbreak{}tree\_\allowbreak{}slide}): every level where the key cells' partition changes is found by bisection over the levels, one connected-component labelling a probe, using the fact that the partition only coarsens as ℓ falls. Every labelling's root count is proved against the tree's own count at that level (\texttt{counts[top − ℓ]}, from the max spanning forest): 0 of 1,102 differ on 44b6\_\allowbreak{}0113de3b, and 0 of 26,454 on 6bba\_\allowbreak{}48816121. On 6bba\_\allowbreak{}48816121, 935 key cells over 100 frames: today's cut (the level with the most components, 626 a frame) holds 931 of them, and only \textbf{50} are alone on their leaf. At each frame's best single level, \textbf{903} are alone. At some level of its own chain, \textbf{925} are alone. Frame 0: the cut is at level 84,314 of 1,135,213, with the 4 key cells on shared leaves. At level 1,045,988 (204 components) all 4 are alone, and they stay separable down to 463,871. So the most-components level is low in the tree: many small pieces, with the cells merged into a few giant blobs. A per-frame or per-chain level separates almost every cell. 44b6\_\allowbreak{}0113de3b has one key cell a frame, so every one is alone at every level (52 of 52). Edges are unchanged at 38/12 and 856/55/7, because nothing in the tracking reads the slide yet. Next is Doug's: the coherence criterion that picks the node per chain, now that the measurement is in.
\end{enumerate}

9c. Answered 23 September, level pairing across frames: “Same residual value”. Cut t and t+1 at the same exact residual r, and count the components of t that overlap exactly one component of t+1 at the frame's drift, and the reverse. \textbf{Measured 23 September} (\texttt{-\allowbreak{}-\allowbreak{}overlap}, \texttt{max\_\allowbreak{}tree\_\allowbreak{}keep\_\allowbreak{}frames} + \texttt{max\_\allowbreak{}tree\_\allowbreak{}overlap}):

\begin{itemize}
\item \textbf{Kept frames.} max\_\allowbreak{}tree keeps the last two frames' codes plus a code→residual table (9 limbs a code), built on the device from the first voxel of each run of equal residuals.
\item \textbf{Threshold.} A residual r maps to each frame's threshold code by one device lane a code: the lane whose value is ≥ r while the value one code below is < r. Unused code slots are zero, and zero is below every admitted residual.
\item \textbf{At each r:}

\begin{itemize}
\item one connected-component labelling a frame;
\item each earlier voxel alive at r whose drifted place is inside the view and alive in the later frame emits the exact pair (earlier root, later root); voxels drifted outside the view emit nothing;
\item a device radix sort, then the distinct pairs, then partner counts a root;
\item “one” is the count of roots with exactly 1 partner, and “mutual” (AND = product) is the count of distinct pairs whose two ends both have exactly 1.
\end{itemize}
\item \textbf{Sampled r values:} each frame's key-partition events from the slide, plus its most-components cut, each mapped to its exact residual. This is a key-driven sample, not every r. Every r would be about 1.1M labellings a frame, and which r to sample is still Doug's.
\item \textbf{6bba\_\allowbreak{}48816121}, 99 pairs, 3,615 cuts:

\begin{itemize}
\item at the earlier frame's cut: 61,679 components, 15,979 overlap exactly one (26\%), and \textbf{6,688 mutual (11\%)};
\item the best mutual over each pair's sampled cuts sums to \textbf{13,783} (139 a pair, against 68 at the cut);
\item frame 0→1 at level 1,045,988: 204 components, 167 one, \textbf{144 mutual}. At the cut, level 84,314: 329 components, 83 one, 29 mutual;
\item the mutual count peaks high in the tree (codes about 1.0 to 1.1M of about 1.1M), at 150 to 200 components a frame, where 70–80\% of the components are mutual. The most-components cut sits low, where 10–30\% are.
\end{itemize}
\item \textbf{44b6\_\allowbreak{}0113de3b}, 53 pairs, 206 cuts, only two sampled levels a frame (the one key cell's appearance and the cut): at the cut 39,507 components, 14,193 one (36\%), 7,962 mutual (20\%); the best mutual sums to 8,166.
\item \textbf{Cost and edges.} The overlap costs about 11 ms a cut: engines went from 125.2 s to 164.3 s on 6bba. Edges are unchanged at 856/55/7 and 38/12, because nothing reads it yet.
\item \textbf{Next is Doug's:} the rule that replaces the most-components cut (for example, the r that maximizes mutual one-to-one, the product of the two verdicts), and whether r is swept over every value or over a data-driven sample.
\end{itemize}

9d. Direction, 23 September (Doug): “make our cell tracker do all the cool stuff organoidtracker does”. OrganoidTracker (vendored at \texttt{vendor\_\allowbreak{}for\_\allowbreak{}ref/\allowbreak{}Organoid\allowbreak{}Tracker}) is the feature list; each feature is rebuilt in our terms: exact integers, verdicts as fields, no weights or cutoffs. The mapping, with nothing built yet:

\begin{itemize}
\item \textbf{Detection.} Theirs: a U-Net predicts centres, then peak calling. Ours: max-tree bodies at the r the overlap picks (items 9, 9c).
\item \textbf{Link score.} Theirs: a link network's probability per candidate pair. Ours: the exact overlap pairs at the drift (9c), with the shared-voxel count as the pair's integer field.
\item \textbf{Division.} Theirs: a division network, then \texttt{pinpoint\_\allowbreak{}divisions}. Ours: in the overlap graph, a component of t whose partners in t+1 are exactly two, each of them with exactly one partner back (1→2).
\item \textbf{Appear and disappear.} Theirs: penalties with a probability floor (0.01) and an edge margin. Ours: 0→1 and 1→0 in the same graph, with view-edge contact read from the body's \texttt{touches} field instead of a margin.
\item \textbf{Global linking.} Theirs: min-cost flow (dpct, Haubold 2016) with five chosen weights. Ours: min-cost flow on integer costs taken from the fields. Where each cost comes from is Doug's.
\item \textbf{Gap bridging.} Theirs: \texttt{bridge\_\allowbreak{}gaps}. Ours: the overlap across t→t+2 at the summed drift, for a 1→0 followed by a 0→1.
\item \textbf{Error probabilities.} Theirs: local marginalization, which enumerates microstates in a 3-step neighbourhood and weights them with Boltzmann factors at temperature 1.5. Ours: the same enumeration with exact counts; a link's field is (valid local solutions that contain it) / (all valid local solutions), kept as an exact pair.
\item \textbf{Manual correction GUI and retraining loop.} Out of scope for the competition.
\item \textbf{I/O, added by Doug: “Ctc geff api queries etc”.}

\begin{itemize}
\item Theirs: \texttt{imaging/\allowbreak{}geff\_\allowbreak{}io.\allowbreak{}py} (\texttt{load\_\allowbreak{}data\_\allowbreak{}file} and \texttt{save\_\allowbreak{}data\_\allowbreak{}file}, with node and edge props), and \texttt{imaging/\allowbreak{}ctc\_\allowbreak{}io.\allowbreak{}py} (reads and writes CTC: the man\_\allowbreak{}track lineage file plus label images).
\item Ours today: \texttt{engine\_\allowbreak{}geff\_\allowbreak{}read} only, used for the answer key. Ours to build: GEFF write of our tracks with node and edge props (the exact fields as props), and CTC read and write.
\end{itemize}
\item \textbf{Query API, added by Doug.}

\begin{itemize}
\item Theirs: \texttt{core/\allowbreak{}links.\allowbreak{}py}.

\begin{itemize}
\item Positions: \texttt{find\_\allowbreak{}futures}, \texttt{find\_\allowbreak{}pasts}, \texttt{find\_\allowbreak{}single\_\allowbreak{}future}, \texttt{find\_\allowbreak{}single\_\allowbreak{}past}.
\item Appear and disappear: \texttt{find\_\allowbreak{}appeared\_\allowbreak{}positions}, \texttt{find\_\allowbreak{}disappeared\_\allowbreak{}positions}.
\item Tracks: \texttt{get\_\allowbreak{}track}, \texttt{find\_\allowbreak{}all\_\allowbreak{}descending\_\allowbreak{}tracks}, \texttt{find\_\allowbreak{}all\_\allowbreak{}tracks\_\allowbreak{}in\_\allowbreak{}same\_\allowbreak{}lineage}, \texttt{find\_\allowbreak{}starting\_\allowbreak{}tracks}, \texttt{find\_\allowbreak{}ending\_\allowbreak{}tracks}, \texttt{find\_\allowbreak{}all\_\allowbreak{}tracks\_\allowbreak{}in\_\allowbreak{}time\_\allowbreak{}point}.
\item Data attached to links and lineages: \texttt{get\_\allowbreak{}link\_\allowbreak{}data}, \texttt{find\_\allowbreak{}all\_\allowbreak{}links\_\allowbreak{}with\_\allowbreak{}data}, \texttt{get\_\allowbreak{}lineage\_\allowbreak{}data}.
\end{itemize}
\item Ours to build: the same queries over our track graph, in one module behind one request struct, with the exact fields as link data.
\end{itemize}
\item \textbf{Order, ruled 23 September: the event census first.} \texttt{-\allowbreak{}-\allowbreak{}overlap} classifies every component at each cut by its partner counts: 1→1 move, 1→2 division, 2→1 merge, 1→0 disappear, 0→1 appear, and the rest as n→m. It is measured against the key's divisions on 44b6 and 6bba.

\begin{itemize}
\item \textbf{Measured 23 September.}

\begin{itemize}
\item Two device passes over the distinct (earlier root, later root) pairs:

\begin{itemize}
\item “backs”: for each root, how many of its partners have exactly one partner in return;
\item a census a root: 0 partners; 1 partner with 1 back (move); 2 partners with 2 backs (1→2, or 2→1 on the later side).
\end{itemize}
\item Mutual counted from the later side equals the earlier side's on every cut, 0 differ on both samples. That is the census's proof.
\end{itemize}
\item \textbf{At the earlier frame's cut:}

\begin{itemize}
\item 6bba\_\allowbreak{}48816121: 6,688 moves, 570 1→2, 590 2→1, 42,853 disappear, 40,996 appear, out of 61,679 components. The key holds 5 divisions over these pairs.
\item 44b6\_\allowbreak{}0113de3b: 7,962 moves, 819 1→2, 789 2→1, 22,545 disappear, 21,517 appear, out of 39,507. The key holds 0 divisions.
\item So at the most-components cut, about 70\% of the components have no partner at the drift: specks that do not persist. Higher in the tree about 20\% do.
\item The key annotates about 9 cells a frame, so census totals cannot be graded against it. Only a per-division verdict can.
\end{itemize}
\item \textbf{Per key division, measured 23 September.} Each division is graded on a chain of verdicts, each implying the one before:

\begin{itemize}
\item present: the parent at t and both daughters at t+1 are above the cut;
\item apart: the daughters are in two distinct components;
\item parted: the parent's component overlaps both of them at the drift, found by a device binary search in the sorted pairs;
\item caught: parted, and the parent's component is 1→2.
\end{itemize}
\item \textbf{6bba\_\allowbreak{}48816121, 5 key divisions, taking each pair's best sampled cut.} At the most-components cut, 0 are parted.

\begin{itemize}
\item 5 are present, 3 apart, 2 parted, 1 caught.
\item t=5→6: caught at level 1,056,677 (214 components, 162 mutual).
\item t=11→12: parted at 3 cuts but never 1→2, because the parent's component has more partners or its partners have more backs.
\item t=57→58: apart at 4 cuts but never parted. The parent's component at the frame's single drift does not reach both daughters.
\item t=16→17 and t=60→61: the daughters are never apart at any cut where the parent is also present.
\end{itemize}
\item \textbf{What that says:}

\begin{itemize}
\item A just-divided pair separates only at levels where the parent is gone, or not at all.
\item One global drift misses daughters that move apart.
\item So a division is found by a per-chain level (row 1's slide), with a per-body drift or climb rather than the frame's.
\end{itemize}
\item \textbf{Next is Doug's:} the rule for the cut; then GEFF and CTC write, and the query API (9d).
\end{itemize}
\end{itemize}

\begin{enumerate}
\item Ruled 23 September (“split it, go”): the slide and the overlap went back into the engine as engine math (the engine book's build plan, item 24). Key cells, divisions and their grading (present, apart, parted, caught) are still to be rebuilt in cell\_\allowbreak{}tracking/src/slide/slide\_\allowbreak{}score.cu from probes and links; that file still holds the old API. They cannot run until score\_\allowbreak{}sample (engine/base/oracle) moves to cell\_\allowbreak{}tracking, since score\_\allowbreak{}sample's frame loop is where they hook in.
\item The driver through tessera, 24 September. Doug: “The engines own runs need to go through tessera”. Done 24 September:

\begin{itemize}
\item track\_\allowbreak{}driver submits \texttt{-\allowbreak{}-\allowbreak{}ingest} and each \texttt{-\allowbreak{}-\allowbreak{}run} part as one job. The signum is BLAKE3(part name, 0, effective cfg). The declaration is the largest sample's lanes × 2: \texttt{engine\_\allowbreak{}source\_\allowbreak{}lanes} (new, shape only) for ingest, the iapx head otherwise.
\item \texttt{-\allowbreak{}-\allowbreak{}override} added. Holding 2 s, sweep 20 ms, idle 5 s. There is no path without tessera.
\item build\_\allowbreak{}driver.sh builds and publishes tessera\_\allowbreak{}daemon beside the driver.
\item Run on a scratch set: ingest declared 838,860,800 with a peak of 5,091,037,184; prove had peaks of 3,962,761,216 and 3,958,566,912. The history is 128 bytes.
\item Finding, not changed (A11 is Doug's rule): a job is reserved at its declaration even when its kept peak is larger. Reserving max(declared, kept peak) is the fix.
\end{itemize}
\item \textbf{The score, 24 September.} Tessera and portability stop here (Doug: “stop getting bogged down in portability”).

\begin{itemize}
\item The real metric is \texttt{maint/\allowbreak{}score\_\allowbreak{}submission.\allowbreak{}py} (edge Jaccard × node count factor). The tracker's own tally (83.1\%) is leaf-to-leaf and is not the score.
\item The current engine reads only \texttt{.\allowbreak{}kcr}. The 25 were re-ingested into \texttt{D:/\allowbreak{}kaggle\_\allowbreak{}project\_\allowbreak{}data/\allowbreak{}biohub\_\allowbreak{}cell\_\allowbreak{}tracking\_\allowbreak{}set\_\allowbreak{}kcr/\allowbreak{}train} (25 of 25 held, 10,207,190,124 bytes, 48.6\% of raw, set root ad40e1d0…f53a39, 757 s); the \texttt{.\allowbreak{}iapx} set is untouched.
\item \textbf{Baseline on the current engine: SCORE 0.209} (object, largest 400; Jaccard 0.210, 961,081 nodes against 880,906). The 22 September engine scored \textbf{0.661} (basin, largest 400; Jaccard 0.670). The tally is identical to 23 September's, so the regression came with the v2 split, not today's work.
\item Cause: the 22 September nodes were peak basins (\texttt{binomial\_\allowbreak{}basins}/\texttt{peak\_\allowbreak{}basins}, steepest ascent over 26 neighbours, membership where the residual is positive), about 300,000 leaves a sample grouped into objects. The v2 engine's nodes are max-tree components at the most-components level, one object a leaf, cells merged into blobs. The basin modules were removed in the split and nothing in the engine does steepest ascent.
\item Built: \texttt{-\allowbreak{}-\allowbreak{}basins} (rule → EngineBuffers → EngineBodiesRequest → MaxTreeObjectsRequest.basins). \texttt{max\_\allowbreak{}tree\_\allowbreak{}basin\_\allowbreak{}ascent\_\allowbreak{}kernel} points every voxel at the highest of itself and its 26 neighbours (key code, then \textasciitilde{}index), and \texttt{max\_\allowbreak{}tree\_\allowbreak{}cc\_\allowbreak{}flatten\_\allowbreak{}kernel} follows each positive voxel to its peak; the peak, mark and body kernels are unchanged. The level cut stays the default until basins score.
\end{itemize}
\end{enumerate}
